\section{Multi-Modal Adversarial Auditor Pipeline Architecture}
\label{appendix:auditor_pipeline}

\begin{figure*}[t]
    \centering
    \includegraphics[width=\linewidth]{auditor-inference.png}
    \caption{
        \textbf{Architecture diagram of the Auditor module.}
    }
    \label{fig:auditor_architecture}
\end{figure*}

This appendix provides a rigorous architectural and mathematical specification
of the Adversarial Image Auditor pipeline. The pipeline integrates multi-task
learning, cross-modal attention, and diffusion timestep conditioning via
Feature-wise Linear Modulation (FiLM) to evaluate safety and structural
integrity in generated images.

% ─────────────────────────────────────────────────────────────────────────────
\subsection{Complete Pipeline Flow and Spatial Dimensions}

The system takes three distinct inputs: an image
$X \in \mathbb{R}^{B \times 3 \times 224 \times 224}$, a tokenized prompt
sequence $T \in \mathbb{R}^{B \times L}$ (where $L$ is the token sequence
length), and a normalized diffusion timestep $t \in \mathbb{R}^{B \times 1}$.
The primary data flow is formalised below.

\begin{enumerate}[leftmargin=*, label=\textbf{\arabic*.}]

\item \textbf{Visual Feature Extraction.}
\begin{equation}
\begin{aligned}
    F_{\text{visual}} = \text{ResNet101}_{\text{backbone}}(X)
    \;\in\; \\ \mathbb{R}^{B \times 2048 \times 7 \times 7} \notag
\end{aligned}
\end{equation}
A global average-pooled vector is also computed: 
\[
\begin{aligned}
f_{\text{global}} = \operatorname{AdaptiveAvgPool2d}(F_{\text{visual}}) \\ \in \mathbb{R}^{B \times 2048} \notag
\end{aligned}
\]

\item \textbf{Textual Encoding.}
\begin{equation}
    f_{\text{text}},\; S_{\text{seq}},\; M_{\text{pad}}
    = \text{BiLSTM}_{\text{encoder}}(T) \notag
\end{equation}
where $f_{\text{text}} \in \mathbb{R}^{B \times 512}$ is the global textual
context, $S_{\text{seq}} \in \mathbb{R}^{B \times L \times 512}$ contains
per-token sequence embeddings (pre-LayerNorm), and
$M_{\text{pad}} \in \{0,1\}^{B \times L}$ is a boolean padding mask.

\item \textbf{Prompt-to-Image Cross-Attention.}
Visual feature maps are projected to the textual dimension, flattened, and
normalised to form query vectors:
\begin{equation}
\begin{aligned}
    Q = \operatorname{LayerNorm}\!\bigl(\operatorname{Conv2d}_{1\times1}(F_{\text{visual}})\bigr) \\
    \;\in\; \mathbb{R}^{B \times 49 \times 512} \notag
\end{aligned}
\end{equation}
\begin{equation}
    K = V = \operatorname{LayerNorm}(S_{\text{seq}})
    \;\in\; \mathbb{R}^{B \times L \times 512} \notag
\end{equation}
Multi-head cross-attention yields prompt-conditioned visual representations:
\begin{equation}
\begin{aligned}
    A_{\text{attended}}
    = \operatorname{MultiHeadAttention}(Q,\,K,\,V;\; \\ \text{mask}=M_{\text{pad}})
    \;\in\; \mathbb{R}^{B \times 49 \times 512} \notag
\end{aligned}
\end{equation}
Spatial averaging gives $f_{\text{attended}} \in \mathbb{R}^{B \times 512}$,
which is used exclusively by the CLIP-style alignment head (Section~B.3).

\end{enumerate}

% ─────────────────────────────────────────────────────────────────────────────
\subsection{Timestep-Aware FiLM Conditioning}

Adversarial noise patterns manifest differently depending on the denoising
timestep $t \in [0,1]$.  A multi-layer perceptron maps $t$ to a shared
temporal embedding:
\begin{equation}
    e_{t} = \operatorname{MLP}_{\text{timestep}}(t)
    \;\in\; \mathbb{R}^{B \times 512} \notag
\end{equation}

This embedding drives two \emph{independent} FiLM projections with separate
learned weights --- one for adversary modulation and one for seam modulation.
Each projection produces scale $\gamma$ and shift $\beta$ coefficients for
Feature-wise Linear Modulation.

\begin{enumerate}[leftmargin=*, label=\textbf{\arabic*.}]

\item \textbf{Relative Adversary Modulation.}
The \emph{relative adversary score} quantifies the continuous perturbation
strength of an image on a $[0,1]$ scale (see Section~B.3).  To condition its
prediction on the timestep, we modulate the global pooled features:
\begin{equation}
\begin{aligned}
     \gamma_{\text{adv}},\;\beta_{\text{adv}}
    = \operatorname{Split}\!\bigl(\operatorname{Linear}_{\text{film,adv}}(e_t)\bigr) \\
    \;\in\; \mathbb{R}^{B \times 2048} \notag
    \end{aligned}
\end{equation}    

\begin{equation}
    f_{\text{global}}^{\text{mod}}
    = (1 + \gamma_{\text{adv}}) \odot f_{\text{global}} + \beta_{\text{adv}} \notag
\end{equation}
Here $\operatorname{Linear}_{\text{film,adv}} : \mathbb{R}^{512} \to
\mathbb{R}^{4096}$, and $\operatorname{Split}$ divides the output equally along
the last dimension.

\item \textbf{Seam Quality Modulation.}
A separate projection conditions the intermediate convolutional features
$F_{\text{seam}} \in \mathbb{R}^{B \times 512 \times 7 \times 7}$ used for
structural artefact detection:
\begin{equation}
\begin{aligned}
    \gamma_{\text{seam}},\;\beta_{\text{seam}}
    = \operatorname{Split}\!\bigl(\operatorname{Linear}_{\text{film,seam}}(e_t)\bigr) \\
    \;\in\; \mathbb{R}^{B \times 512} \notag
\end{aligned}
\end{equation}
\begin{equation}
    F_{\text{seam}}^{\text{mod}}
    = \bigl(1 + \gamma_{\text{seam}}^{\phantom{A}}\bigr) \odot F_{\text{seam}}
      + \beta_{\text{seam}} \notag
\end{equation}
where $\gamma_{\text{seam}}$ and $\beta_{\text{seam}}$ are broadcast over the
spatial dimensions.  Note that $\operatorname{Linear}_{\text{film,seam}}$ maps
$\mathbb{R}^{512} \to \mathbb{R}^{1024}$ and has no shared weights with
$\operatorname{Linear}_{\text{film,adv}}$.

\end{enumerate}

% ─────────────────────────────────────────────────────────────────────────────
\subsection{Specialised Multi-Task Auditor Heads}

The auditor decouples distinct tasks into dedicated prediction heads to ensure
stable gradient flow and to protect alignment encodings from direct safety-label
contamination.

\begin{itemize}[leftmargin=*]

\item \textbf{Binary Adversarial Head.}
Computes $P(\text{Adversarial} \mid X)$ directly from the \emph{unmodulated}
spatial features $F_{\text{visual}}$ via a $1{\times}1$ convolution, adaptive
global average pooling, and Sigmoid activation.  This head intentionally
bypasses FiLM conditioning so that the binary decision is not conflated with
timestep dynamics.

\item \textbf{Safety Category Head.}
Projects $F_{\text{visual}}$ to $C=3$ channels and applies cost-sensitive
cross-entropy with class weights $w = [1.0,\,5.0,\,2.0]$.  The elevated weight
for class~1 (adversarial) reflects its underrepresentation in the training
distribution; the moderate weight for class~2 (borderline) penalises
ambiguous predictions more than the majority safe class.

\item \textbf{Relative Adversary Score Head.}
Regresses the continuous perturbation strength $s \in [0,1]$ from the
FiLM-modulated vector $f_{\text{global}}^{\text{mod}}$ using a three-layer MLP
with dropout.  Unlike the binary head, this head is conditioned on the timestep
embedding to capture how perturbation magnitude varies across denoising stages.

\item \textbf{Seam Quality Assessment Head.}
Evaluates inpainting seams and composite boundaries by applying a regression
head to $F_{\text{seam}}^{\text{mod}}$, yielding a localised artefact score.

\item \textbf{CLIP-style Faithfulness Alignment Head.}
Projects $f_{\text{attended}}$ and $f_{\text{text}}$ independently into a
shared latent space $\mathbb{R}^{256}$ using separate two-layer MLPs, then
optimises image--prompt alignment via symmetric InfoNCE loss with a learnable
log-temperature $\tau$.  Both the cross-attended visual vector and the raw
global text vector $f_{\text{text}}$ are fed to this head to preserve low-level
lexical grounding alongside spatially-attended semantics.

\end{itemize}

% ─────────────────────────────────────────────────────────────────────────────
\subsection{GradCAM Explainability Pipeline}

For visual explainability, target activations are extracted from the final
residual convolutional layer of the ResNet101 backbone.  The localised risk
heatmap $M_{\text{risk}} \in \mathbb{R}^{7 \times 7}$ for safety class $c$ is:
\begin{equation}
    M_{\text{risk}}
    = \operatorname{ReLU}\!\left(\sum_{k} \alpha_k^c \cdot F_{\text{visual},k}\right) \notag
\end{equation}
where the importance weights $\alpha_k^c$ are obtained by global average
pooling of the backpropagated gradients:
\begin{equation}
    \alpha_k^c
    = \operatorname{AdaptiveAvgPool2d}\!\left(
        \frac{\partial\, y^c}{\partial\, F_{\text{visual},k}}
      \right) \notag
\end{equation}
The resulting maps are bilinearly upsampled to $224{\times}224$ and
alpha-blended onto the source image using custom contour overlays to localise
flagged regions.  Because GradCAM operates directly on $F_{\text{visual}}$, it
provides explanations that are independent of FiLM conditioning, giving an
unbiased view of which image regions drive the safety classification.

